\hypertarget{user-guide---rnra-data-processing-and-error-management}{%
\section{USER GUIDE - RNRA data processing and error
management}\label{user-guide---rnra-data-processing-and-error-management}}

\hypertarget{table-of-contents}{%
\subsection{Table of Contents}\label{table-of-contents}}

\begin{itemize}
\tightlist
\item
  \protect\hyperlink{1-introduction}{1. Introduction}
\item
  \protect\hyperlink{2-scientific-background}{2. Scientific Background}
\item
  \protect\hyperlink{3-installation}{3. Installation}
\item
  \protect\hyperlink{4-interface-description}{4. Interface Description}
\item
  \protect\hyperlink{5-processing-workflow}{5. Processing Workflow}
\item
  \protect\hyperlink{6-output-files}{6. Output Files}
\item
  \protect\hyperlink{7-methodological-details}{7. Methodological
  Details}
\item
  \protect\hyperlink{8-Code-Architecture}{8. Code Architecture}
\item
  \protect\hyperlink{9-limitations}{9. Limitations}
\item
  \protect\hyperlink{10-references}{10. References}
\end{itemize}

\hypertarget{introduction}{%
\subsection{1. Introduction}\label{introduction}}

\begin{itemize}
\tightlist
\item
  objectif du logiciel.
\item
  publique cible
\end{itemize}

\hypertarget{scientific-principle}{%
\subsection{2. scientific principle}\label{scientific-principle}}

Explication : - RNRA - Methode ANOVA - Hypothese statistique - Test de
normalité et homogénité

\hypertarget{detailed-installation}{%
\subsection{3. detailed installation}\label{detailed-installation}}

Procedure etape par etape

Here is a cleaned‑up, up‑to‑date English markdown draft for the key
``user guide'' parts, aligned with your current GUI (Conversion /
Calibration / Processing). You can paste these blocks into
\texttt{USER\_GUIDE.md} and adapt the numbering if needed.

\begin{center}\rule{0.5\linewidth}{0.5pt}\end{center}

\hypertarget{interface-description}{%
\subsection{4. Interface Description}\label{interface-description}}

\hypertarget{conversion-window}{%
\subsubsection{4.1 Conversion Window}\label{conversion-window}}

This window performs the batch conversion of raw \texttt{.mpa}
acquisition files into structured \texttt{.txt} spectra that can be used
in the calibration and processing steps. The conversion routine reads
the acquisition metadata, computes the dead time correction factor, and
exports one text file per ADC channel with a standardized format.

\hypertarget{layout-and-controls}{%
\paragraph{Layout and controls}\label{layout-and-controls}}

The \textbf{Conversion} tab contains the following elements:

\begin{itemize}
\tightlist
\item
  \textbf{Configuration file (Excel)}

  \begin{itemize}
  \tightlist
  \item
    Read‑only text field displaying the path to the global configuration
    Excel file.\\
  \item
    \textbf{Browse} button to select the configuration file
    (e.g.~\texttt{config.xlsx}).
  \end{itemize}
\item
  \textbf{Day root (sample name root)}

  \begin{itemize}
  \tightlist
  \item
    Text entry labelled e.g.~\texttt{day\ racine\ samplename} where the
    user enters the day or sample root, such as \texttt{250317}.\\
  \item
    This value is used to filter the rows corresponding to a given
    measurement day in the configuration file.
  \end{itemize}
\item
  \textbf{Output folder for .txt files}

  \begin{itemize}
  \tightlist
  \item
    Text entry labelled \texttt{Folder\ exit\ for\ .txt} displaying the
    path where converted \texttt{.txt} files will be written.\\
  \item
    \textbf{Browse} button to choose or create this directory.
  \end{itemize}
\item
  \textbf{Conversion controls}

  \begin{itemize}
  \tightlist
  \item
    \textbf{convert .mpa → .txt} button to start the conversion for the
    selected day.\\
  \item
    Progress bar indicating the current file index and total number of
    files being processed.\\
  \item
    Status label showing messages such as ``Waiting'', ``Conversion in
    progress\ldots{}'', ``Finished'', or ``Finished with errors''.
  \end{itemize}
\item
  \textbf{Log window}

  \begin{itemize}
  \tightlist
  \item
    Multi‑line text area showing detailed messages, including which
    folders are used, which files are converted, and any encountered
    errors.
  \end{itemize}
\end{itemize}

\hypertarget{user-actions}{%
\paragraph{User actions}\label{user-actions}}

\begin{enumerate}
\def\labelenumi{\arabic{enumi}.}
\tightlist
\item
  \textbf{Load the configuration file}

  \begin{itemize}
  \tightlist
  \item
    Click \textbf{Browse} in the ``config Excel'' section.\\
  \item
    Select the Excel configuration file used to describe your
    measurements (sample names, folders, first/last file numbers,
    etc.).\\
  \item
    The selected path appears in the read‑only entry.
  \end{itemize}
\item
  \textbf{Define the day root}

  \begin{itemize}
  \tightlist
  \item
    In the \textbf{day root} entry, type the day/sample root
    corresponding to the measurement series you want to convert (for
    example \texttt{250317}).\\
  \item
    Internally, the software looks for rows in the configuration file
    whose sample name starts with this root and retrieves the
    corresponding \texttt{mpafolder} path.
  \end{itemize}
\item
  \textbf{Choose the output folder for \texttt{.txt} files}

  \begin{itemize}
  \tightlist
  \item
    In \textbf{Folder exit for .txt}, click \textbf{Browse} and select a
    base directory in which the converted files will be stored.\\
  \item
    The program will create (if necessary) a subfolder named after the
    day root (e.g.~\texttt{\ldots{}/250317/}) and place all
    \texttt{.txt} outputs there.
  \end{itemize}
\item
  \textbf{Launch the conversion}

  \begin{itemize}
  \tightlist
  \item
    Click \textbf{convert .mpa → .txt}.\\
  \item
    The software scans the \texttt{mpafolder} corresponding to the
    selected day in the configuration file, then converts all
    \texttt{.mpa} files found in this folder.\\
  \item
    Progress is displayed in the progress bar and status label; detailed
    messages appear in the log.
  \end{itemize}
\item
  \textbf{Check for errors}

  \begin{itemize}
  \tightlist
  \item
    At the end of the process, if errors occurred (missing folder,
    unreadable \texttt{.mpa} file, etc.), they are listed in the log
    area.\\
  \item
    The status label indicates whether the conversion finished
    successfully or with errors.
  \end{itemize}
\end{enumerate}

\hypertarget{output-text-file-structure}{%
\paragraph{Output text file
structure}\label{output-text-file-structure}}

For each \texttt{.mpa} file and for each ADC channel, one \texttt{.txt}
file is generated.

\begin{itemize}
\item
  \textbf{File naming}

  \begin{itemize}
  \tightlist
  \item
    General convention:\\
    \texttt{filename\_ADCname.txt}~\\
  \item
    \texttt{filename}: base name of the original \texttt{.mpa} file
    (without extension).\\
  \item
    \texttt{ADCname}: ADC channel identifier (\texttt{ADC0},
    \texttt{ADC1}, etc.).
  \end{itemize}

  Example:

  \begin{itemize}
  \tightlist
  \item
    Input file: \texttt{sample01.mpa} with two ADC channels
    (\texttt{ADC1}, \texttt{ADC2})\\
  \item
    Output files:

    \begin{itemize}
    \tightlist
    \item
      \texttt{sample01\_ADC1.txt}~\\
    \item
      \texttt{sample01\_ADC2.txt}
    \end{itemize}
  \end{itemize}
\item
  \textbf{File content}

  \begin{enumerate}
  \def\labelenumi{\arabic{enumi}.}
  \tightlist
  \item
    \textbf{Header line}

    \begin{itemize}
    \tightlist
    \item
      First line, typically containing the dead time factor extracted
      from the \texttt{.mpa} metadata, such as:\\
      \texttt{Dead\ time\ factor\ =\ 1.023}
    \end{itemize}
  \item
    \textbf{Data lines} (one measurement per line)

    \begin{itemize}
    \tightlist
    \item
      \texttt{Channel}: integer ADC channel index.\\
    \item
      \texttt{Count}: measured counts in this channel.
    \end{itemize}
  \end{enumerate}
\end{itemize}

All exported text files are formatted to be directly readable in Python
(via \texttt{numpy.loadtxt}, \texttt{pandas.read\_csv}) or in other
scientific data processing tools.

\begin{center}\rule{0.5\linewidth}{0.5pt}\end{center}

\hypertarget{calibration-window}{%
\subsubsection{4.2 Calibration Window}\label{calibration-window}}

The \textbf{Calibration} tab is used to build an energy calibration
curve \(E = a \cdot C + b\) from one or several spectra in \texttt{.txt}
format generated by the conversion step. The user interactively selects
peaks on a displayed spectrum, fits a Gaussian shape to each peak to
extract the centroid channel, associates each centroid with a known
reference energy, then performs a weighted linear regression.

\hypertarget{layout-and-controls-1}{%
\paragraph{Layout and controls}\label{layout-and-controls-1}}

The \textbf{Calibration} window is organised into three areas:

\begin{itemize}
\item
  \textbf{Left panel -- File and peak management}

  \begin{itemize}
  \tightlist
  \item
    \textbf{``Load spectrum''} button

    \begin{itemize}
    \tightlist
    \item
      Opens a file dialog to select a spectrum \texttt{.txt} file (one
      ADC channel) produced by the conversion window.\\
    \item
      After loading, the filename and dead time factor are displayed.
    \end{itemize}
  \item
    \textbf{File information labels}

    \begin{itemize}
    \tightlist
    \item
      Label showing the currently loaded spectrum file name.\\
    \item
      Label showing the associated dead time factor if present in the
      header.
    \end{itemize}
  \item
    \textbf{Peak selection area}

    \begin{itemize}
    \tightlist
    \item
      Instruction label (``Click on the plot to select a peak'').\\
    \item
      Label showing the currently selected interval in channels (center
      and width).\\
    \item
      Entry field for the \textbf{reference energy (keV)} corresponding
      to the selected peak.\\
    \item
      \textbf{Fit Gaussian} button to launch the peak fit for the
      current selection.\\
    \item
      \textbf{Reset selection} button to clear the current selection.
    \end{itemize}
  \item
    \textbf{Peak list (Treeview)}

    \begin{itemize}
    \tightlist
    \item
      Table with one row per identified calibration peak.\\
    \item
      Columns:

      \begin{itemize}
      \tightlist
      \item
        Index (internal peak number)\\
      \item
        \textbf{Channel} (fitted centroid)\\
      \item
        \textbf{Energy (keV)} (reference value entered by the user)
      \end{itemize}
    \end{itemize}
  \item
    \textbf{Peak management buttons}

    \begin{itemize}
    \tightlist
    \item
      \textbf{Delete selected peak}: removes the highlighted peak from
      the list and from the internal data.\\
    \item
      \textbf{Clear all peaks}: removes all peaks and resets the
      calibration state.
    \end{itemize}
  \item
    \textbf{Calibration controls}

    \begin{itemize}
    \tightlist
    \item
      \textbf{Run calibration} button (enabled when at least two peaks
      are defined).\\
    \item
      Label showing the calibration result or a message such as
      ``Minimum 2 peaks required''.
    \end{itemize}
  \end{itemize}
\item
  \textbf{Center panel -- Spectrum and fits}

  \begin{itemize}
  \tightlist
  \item
    Matplotlib figure displaying the loaded spectrum (counts
    vs.~channel).\\
  \item
    Tools for zoom, pan, etc. (Matplotlib toolbar).\\
  \item
    An interactive \textbf{SpanSelector} allowing the user to select a
    channel interval by clicking and dragging on the plot.
  \item
    After fitting, vertical lines and Gaussian curves are drawn to show
    the fitted peaks.
  \end{itemize}
\item
  \textbf{Right panel -- Detailed results and export}

  \begin{itemize}
  \tightlist
  \item
    Text area listing:

    \begin{itemize}
    \tightlist
    \item
      General information about the loaded spectrum.\\
    \item
      Details of each peak fit (centroid, error, sigma, amplitude).\\
    \item
      Calibration results (slope, intercept, goodness of fit, relative
      RMS).\\
    \end{itemize}
  \item
    \textbf{Export results} button to save the list of peaks and the
    calibration summary in a text file.
  \end{itemize}
\end{itemize}

\hypertarget{user-actions-1}{%
\paragraph{User actions}\label{user-actions-1}}

\begin{enumerate}
\def\labelenumi{\arabic{enumi}.}
\tightlist
\item
  \textbf{Load a spectrum}

  \begin{itemize}
  \tightlist
  \item
    Click \textbf{Load spectrum} and select a \texttt{.txt} file
    produced by the conversion window.\\
  \item
    The program reads the header and data, displays the spectrum, and
    shows the dead time factor if available.\\
  \item
    The plot is updated to show counts versus channel index.
  \end{itemize}
\item
  \textbf{Select a peak}

  \begin{itemize}
  \tightlist
  \item
    Option A: Click and drag on the spectrum to select a peak region
    (SpanSelector).

    \begin{itemize}
    \tightlist
    \item
      The selected center and approximate half‑width (in channels) are
      computed automatically.\\
    \end{itemize}
  \item
    Option B: Simply click on the spectrum to set an approximate center;
    an internal rule computes a suitable half‑width depending on the
    channel.\\
  \item
    The current selection is displayed as
    \texttt{center\ =\ \ldots{},\ width\ =\ \ldots{}} in the
    corresponding label.
  \end{itemize}
\item
  \textbf{Associate a reference energy}

  \begin{itemize}
  \tightlist
  \item
    In the \textbf{Energy (keV)} entry, type the known energy of the
    selected peak (e.g.~a resonance or gamma line).\\
  \item
    This value will be used as the calibration reference for this peak.
  \end{itemize}
\item
  \textbf{Fit a Gaussian peak}

  \begin{itemize}
  \tightlist
  \item
    Click \textbf{Fit Gaussian}.\\
  \item
    The program extracts the data around the selected center, performs a
    Gaussian fit, and returns:

    \begin{itemize}
    \tightlist
    \item
      Centroid \(C\) (fitted channel),\\
    \item
      Sigma (peak width),\\
    \item
      Amplitude,\\
    \item
      Uncertainty on the centroid.\\
    \end{itemize}
  \item
    The peak is added to the Treeview with its centroid and reference
    energy, and details are appended to the text area.
  \end{itemize}
\item
  \textbf{Repeat for all calibration peaks}

  \begin{itemize}
  \tightlist
  \item
    Repeat steps 2--4 for each reference line you want to include in the
    calibration.\\
  \item
    Once at least two peaks are defined, the \textbf{Run calibration}
    button is enabled.
  \end{itemize}
\item
  \textbf{Run the linear calibration}

  \begin{itemize}
  \tightlist
  \item
    Click \textbf{Run calibration}.\\
  \item
    The software performs a (weighted) linear regression
    \(E = a \cdot C + b\) using the centroids and their uncertainties.\\
  \item
    It computes:

    \begin{itemize}
    \tightlist
    \item
      Slope \(a\) and intercept \(b\),\\
    \item
      Uncertainties on \(a\) and \(b\),\\
    \item
      Coefficient of determination \(R^2\),\\
    \item
      Global RMS error and relative RMS (in \%).\\
    \end{itemize}
  \item
    A summary string is displayed in the calibration result label and in
    the text area.
  \end{itemize}
\item
  \textbf{Update the configuration file (if applicable)}

  \begin{itemize}
  \tightlist
  \item
    If the application is used together with the Conversion and
    Processing tabs, and a global configuration Excel file is defined,
    the calibration results can be written back into the configuration
    file for the current group/day.\\
  \item
    The corresponding columns (e.g.~\texttt{slope}, \texttt{intercept},
    \texttt{errorcalib}) are updated for the matching rows, so the
    processing loop can use them automatically.
  \end{itemize}
\item
  \textbf{Export calibration results (optional)}

  \begin{itemize}
  \tightlist
  \item
    Click \textbf{Export results} to save:

    \begin{itemize}
    \tightlist
    \item
      The list of peaks (index, centroid, sigma, amplitude, energy,
      centroid error),\\
    \item
      The linear calibration summary (slope, intercept, uncertainties,
      \(R^2\), relative RMS).\\
    \end{itemize}
  \item
    The export is written to a user‑selected text file.
  \end{itemize}
\end{enumerate}

\hypertarget{output-and-stored-parameters}{%
\paragraph{Output and stored
parameters}\label{output-and-stored-parameters}}

\begin{itemize}
\tightlist
\item
  The main \textbf{numerical result} is a linear calibration
  \(E = a \cdot C + b\) with associated uncertainties and quality
  indicators.
\item
  These parameters are either stored in memory (for immediate use) or
  written into the configuration Excel file to be reused by the
  processing loop.
\end{itemize}

\begin{center}\rule{0.5\linewidth}{0.5pt}\end{center}

\hypertarget{processing-loop-peak-removal-sigmoid-fit}{%
\subsubsection{4.3 Processing (Loop, Peak Removal, Sigmoid
Fit)}\label{processing-loop-peak-removal-sigmoid-fit}}

The \textbf{Processing} window controls the full data‑processing
pipeline, from ROI integration to optional peak removal and final
sigmoid fitting of the excitation profiles. It uses the configuration
Excel file and the calibration parameters to convert profiles from
channel space to energy space and to generate processed outputs for
further analysis.

\hypertarget{layout-and-controls-2}{%
\paragraph{Layout and controls}\label{layout-and-controls-2}}

The \textbf{Processing / Loop} tab is structured as follows:

\begin{itemize}
\item
  \textbf{Configuration and ROI (left panel)}

  \begin{itemize}
  \tightlist
  \item
    \textbf{Configuration file} section

    \begin{itemize}
    \tightlist
    \item
      Read‑only entry showing the path to the configuration Excel
      file.\\
    \item
      \textbf{Load Excel configuration file} button.
    \end{itemize}
  \item
    \textbf{ROI selection} section

    \begin{itemize}
    \tightlist
    \item
      Entry showing or allowing manual edit of the ROI in energy
      (e.g.~\texttt{Emin–Emax\ keV}).\\
    \item
      \textbf{Choose ROI on a spectrum} button to open an interactive
      window where the user selects the ROI directly on a reference
      spectrum (energy vs counts).
    \end{itemize}
  \item
    \textbf{Results saving options}

    \begin{itemize}
    \tightlist
    \item
      Check box \textbf{Save results} to enable/disable saving of
      processed profiles.\\
    \item
      Output folder entry and \textbf{Browse} button (enabled only if
      ``Save results'' is checked).
    \end{itemize}
  \item
    \textbf{Loop execution}

    \begin{itemize}
    \tightlist
    \item
      \textbf{Process files} (or similar label) button to start the loop
      over all scenarios described in the configuration file.
    \end{itemize}
  \end{itemize}
\item
  \textbf{Carbon build‑up peak removal}

  \begin{itemize}
  \tightlist
  \item
    Input fields:

    \begin{itemize}
    \tightlist
    \item
      \textbf{Energy center (keV)}: resonance energy from which the
      build‑up peak is removed (default \textasciitilde6385 keV for the
      15N(H, α, γ)12C resonance).\\
    \item
      \textbf{Search window (keV)}: total energy window around the
      center where the peak search is performed.\\
    \item
      \textbf{Removal half‑width (keV)}: half‑width of the interval to
      be removed once a peak is detected.\\
    \end{itemize}
  \item
    \textbf{Remove build‑up peak} button to apply the removal to all
    profiles previously calculated by the loop.
  \end{itemize}
\item
  \textbf{Sigmoid fitting}

  \begin{itemize}
  \tightlist
  \item
    \textbf{Run sigmoid fit on output profiles} button.\\
  \item
    Uses either the raw loop output or the ``filtered'' folder (after
    peak removal) as input.\\
  \item
    Produces sigmoid fits and associated parameters for each profile.
  \end{itemize}
\item
  \textbf{Visualization options}

  \begin{itemize}
  \tightlist
  \item
    \textbf{Display output profile} button to open a file dialog and
    select an output \texttt{.xlsx} or \texttt{.png} file.\\
  \item
    The selected file is displayed in the Matplotlib figure on the right
    (profiles or images).
  \end{itemize}
\item
  \textbf{Right panel -- Plot and detailed log}

  \begin{itemize}
  \tightlist
  \item
    Matplotlib figure with axes labelled ``Energy (keV)'' and ``NC'' or
    similar.\\
  \item
    Toolbar for zoom/pan.\\
  \item
    Text area showing detailed log messages (loop progress, errors, fit
    results, etc.).
  \end{itemize}
\end{itemize}

\hypertarget{user-actions-2}{%
\paragraph{User actions}\label{user-actions-2}}

\hypertarget{a.-loop-over-roi}{%
\subparagraph{A. Loop over ROI}\label{a.-loop-over-roi}}

\begin{enumerate}
\def\labelenumi{\arabic{enumi}.}
\tightlist
\item
  \textbf{Load the configuration file}

  \begin{itemize}
  \tightlist
  \item
    In the \textbf{Configuration file} section, click \textbf{Load Excel
    configuration file} and select the same Excel file used
    previously.\\
  \item
    The program checks that all required columns are present (sample
    name, tension folder, data folder, first file, last file, ADC,
    slope, intercept, calibration error, etc.).
  \end{itemize}
\item
  \textbf{Define the ROI in energy}

  \begin{itemize}
  \tightlist
  \item
    Option A -- \textbf{Manual}: type the ROI in keV in the
    corresponding entry (e.g.~\texttt{6400–6600\ keV}).\\
  \item
    Option B -- \textbf{Interactive}: click \textbf{Choose ROI on a
    spectrum} to open a dedicated window:

    \begin{itemize}
    \tightlist
    \item
      A reference spectrum (converted to energy using the calibration)
      is displayed.\\
    \item
      You can click‑and‑drag to select the energy interval; the selected
      \texttt{Emin} and \texttt{Emax} are shown and can be refined
      numerically.\\
    \item
      When validated, the ROI is stored both in energy and in
      corresponding channel limits.
    \end{itemize}
  \end{itemize}
\item
  \textbf{Choose whether to save results}

  \begin{itemize}
  \tightlist
  \item
    If you want to keep all intermediate and final results, check
    \textbf{Save results} and choose an output folder.\\
  \item
    If not checked, the program uses an internal temporary directory for
    the session.
  \end{itemize}
\item
  \textbf{Run the processing loop}

  \begin{itemize}
  \tightlist
  \item
    Click \textbf{Process files}.\\
  \item
    For each scenario (row) in the configuration file, the program:

    \begin{itemize}
    \tightlist
    \item
      Reads the relevant \texttt{.txt} spectra.\\
    \item
      Integrates counts within the selected ROI.\\
    \item
      Normalizes by the integrated charge (or other
      intensity/normalization parameter).\\
    \item
      Builds an excitation profile as a function of beam energy.\\
    \end{itemize}
  \item
    Results are written to the output folder (one file per scenario,
    plus summary files if implemented) and messages are displayed in the
    log field.
  \end{itemize}
\end{enumerate}

\hypertarget{b.-peak-removal-optional}{%
\subparagraph{B. Peak removal
(optional)}\label{b.-peak-removal-optional}}

\begin{enumerate}
\def\labelenumi{\arabic{enumi}.}
\tightlist
\item
  \textbf{Set the build‑up energy and windows}

  \begin{itemize}
  \tightlist
  \item
    In \textbf{Energy center (keV)}, keep the default resonance energy
    (e.g.~6385 keV) or replace it with the value relevant for your
    experiment.\\
  \item
    Adjust the \textbf{search window} and \textbf{removal half‑width} if
    needed (default values are provided in the GUI).
  \end{itemize}
\item
  \textbf{Apply the removal to all profiles}

  \begin{itemize}
  \tightlist
  \item
    Ensure the loop has been run at least once (raw profiles exist).\\
  \item
    Click \textbf{Remove build‑up peak}.\\
  \item
    The program scans the output profiles in the selected folder,
    identifies peaks around the given energy, and removes the
    corresponding points in the specified interval.\\
  \item
    Filtered profiles are saved into a dedicated subfolder
    (e.g.~\texttt{filtered}) along with optional diagnostic plots.
  \end{itemize}
\end{enumerate}

\hypertarget{c.-sigmoid-fit}{%
\subparagraph{C. Sigmoid fit}\label{c.-sigmoid-fit}}

\begin{enumerate}
\def\labelenumi{\arabic{enumi}.}
\tightlist
\item
  \textbf{Select the input folder for fitting}

  \begin{itemize}
  \tightlist
  \item
    If peak removal was performed, the fitting step uses the
    \textbf{filtered} folder containing cleaned profiles.\\
  \item
    Otherwise, it uses the raw output folder of the loop.
  \end{itemize}
\item
  \textbf{Run the sigmoid fits}

  \begin{itemize}
  \tightlist
  \item
    Click \textbf{Run sigmoid fit on output profiles}.\\
  \item
    For each excitation profile, the program fits a sigmoid curve and
    extracts parameters such as:

    \begin{itemize}
    \tightlist
    \item
      Low and high plateau values,\\
    \item
      Midpoint energy and width,\\
    \item
      Difference between high and low plateau (used to group profiles or
      compare conditions).\\
    \end{itemize}
  \item
    Fit results and diagnostics are saved to the output directory.
  \end{itemize}
\end{enumerate}

\hypertarget{d.-visualization}{%
\subparagraph{D. Visualization}\label{d.-visualization}}

\begin{enumerate}
\def\labelenumi{\arabic{enumi}.}
\tightlist
\item
  \textbf{Open a processed profile or image}

  \begin{itemize}
  \tightlist
  \item
    Click \textbf{Display output profile}.\\
  \item
    Choose either:

    \begin{itemize}
    \tightlist
    \item
      An Excel file (\texttt{.xlsx}/\texttt{.xls}) containing a profile
      (columns such as \texttt{Energy\ keV}, \texttt{NC}, optionally
      \texttt{uncertainties}), or\\
    \item
      An image file (\texttt{.png}, \texttt{.jpg}, \texttt{.jpeg})
      generated by the processing pipeline.
    \end{itemize}
  \end{itemize}
\item
  \textbf{Inspect the result}

  \begin{itemize}
  \tightlist
  \item
    If an Excel file is selected, the program plots the profile,
    optionally with error bars, and can overlay the fitted curve if
    present.\\
  \item
    If an image file is selected, it is displayed directly in the
    plotting area (axes are hidden for a clean display).
  \end{itemize}
\end{enumerate}

\hypertarget{output-file-naming-and-content-processing}{%
\paragraph{Output file naming and content
(processing)}\label{output-file-naming-and-content-processing}}

The exact naming conventions may evolve, but in general the processing
step produces:

\begin{itemize}
\tightlist
\item
  \textbf{Loop output files}

  \begin{itemize}
  \tightlist
  \item
    Typically one file per scenario/group, containing energy, normalized
    counts, and optional uncertainties.
  \end{itemize}
\item
  \textbf{Filtered profiles} (after peak removal)

  \begin{itemize}
  \tightlist
  \item
    Saved in a \texttt{filtered} subfolder, with the same or similar
    naming scheme as the raw outputs.
  \end{itemize}
\item
  \textbf{Sigmoid fit outputs}

  \begin{itemize}
  \tightlist
  \item
    Files or images summarizing the fit parameters and the quality of
    the fit for each profile.
  \end{itemize}
\end{itemize}

These outputs are then used for further analysis, comparison between
samples, and statistical processing (including ANOVA) as described in
later sections of the guide.

\begin{center}\rule{0.5\linewidth}{0.5pt}\end{center}

\hypertarget{processing-workflow-overview}{%
\subsection{5. Processing Workflow
(Overview)}\label{processing-workflow-overview}}

This section summarises the recommended end‑to‑end workflow using the
three main windows.

\begin{enumerate}
\def\labelenumi{\arabic{enumi}.}
\tightlist
\item
  \textbf{Configure and convert}

  \begin{itemize}
  \tightlist
  \item
    Prepare or update the Excel configuration file describing all
    measurements (paths, file numbers, ADC, calibration parameters).\\
  \item
    Use the \textbf{Conversion} window to convert all \texttt{.mpa}
    files for a given day/root into \texttt{.txt} spectra.
  \end{itemize}
\item
  \textbf{Calibrate energy}

  \begin{itemize}
  \tightlist
  \item
    In the \textbf{Calibration} window, load representative spectra and
    identify several reference peaks.\\
  \item
    Fit Gaussian peaks, assign reference energies, and run the linear
    calibration to obtain \(E = a \cdot C + b\).\\
  \item
    Store the calibration coefficients in the configuration file.
  \end{itemize}
\item
  \textbf{Process profiles}

  \begin{itemize}
  \tightlist
  \item
    In the \textbf{Processing} window, load the configuration file.\\
  \item
    Define the ROI in energy and run the loop to build excitation
    profiles for all scenarios.\\
  \item
    Optionally, remove the carbon build‑up peak around a given energy
    and re‑use the cleaned profiles.
  \end{itemize}
\item
  \textbf{Fit and analyse}

  \begin{itemize}
  \tightlist
  \item
    Perform sigmoid fits on the final excitation profiles to extract
    plateau differences and key parameters.\\
  \item
    Use the exported Excel/text files and images for further analysis,
    comparison between conditions, and statistical evaluation (ANOVA,
    uncertainty propagation, etc.).
  \end{itemize}
\end{enumerate}

\begin{center}\rule{0.5\linewidth}{0.5pt}\end{center}

\hypertarget{output-files-overview}{%
\subsection{6. Output Files (Overview)}\label{output-files-overview}}

The application generates several categories of output files at
different stages.

\begin{itemize}
\tightlist
\item
  \textbf{Converted spectra (\texttt{.txt})}

  \begin{itemize}
  \tightlist
  \item
    Location: subfolder of the chosen output directory, named after the
    day root (e.g.~\texttt{\ldots{}/250317/}).\\
  \item
    One file per \texttt{.mpa} and per ADC channel, with dead time
    factor in the header and \texttt{(channel,\ count)} pairs in the
    body.
  \end{itemize}
\item
  \textbf{Calibration exports}

  \begin{itemize}
  \tightlist
  \item
    Optional text files containing the list of calibration peaks and the
    fitted linear calibration parameters.\\
  \item
    Calibration coefficients may also be written back into the
    configuration Excel file (columns such as slope, intercept,
    calibration error).
  \end{itemize}
\item
  \textbf{Loop / ROI integration outputs}

  \begin{itemize}
  \tightlist
  \item
    Files (often Excel) containing excitation profiles: energy,
    normalized counts, and uncertainties.\\
  \item
    Often organized per sample or per scenario as defined in the
    configuration file.
  \end{itemize}
\item
  \textbf{Filtered profiles (after peak removal)}

  \begin{itemize}
  \tightlist
  \item
    Profiles where points in the build‑up peak region have been removed,
    stored in a dedicated subfolder.
  \end{itemize}
\item
  \textbf{Sigmoid fit results and images}

  \begin{itemize}
  \tightlist
  \item
    Numerical summaries of fit parameters (plateau values, midpoint,
    width, residuals, etc.).\\
  \item
    Plots (\texttt{.png}/\texttt{.jpg}) showing data points and fitted
    curves for visual inspection.
  \end{itemize}
\end{itemize}

These outputs can be directly used in subsequent analysis scripts
(Python, R, etc.) and in statistical post‑processing such as ANOVA and
uncertainty propagation.

\hypertarget{methodological-details}{%
\subsection{7. Methodological Details}\label{methodological-details}}

\hypertarget{conversion-and-dead-time-correction}{%
\subsubsection{7.1 Conversion and Dead Time
Correction}\label{conversion-and-dead-time-correction}}

\hypertarget{input-data}{%
\paragraph{Input data}\label{input-data}}

The conversion step starts from raw \texttt{.mpa} files produced by the
acquisition system. These files contain several structured sections
associated with the different detector channels, including metadata such
as live time and real time for each ADC block.

A typical detector section may contain information of the form:

\begin{verbatim}
[ADC1]
livetime = ...
realtime = ...
\end{verbatim}

During conversion, the software parses the file to identify each ADC
section, extracts the corresponding acquisition metadata, and exports
the spectrum as a structured text file containing a header and
channel-count pairs.

\hypertarget{mathematical-transformation}{%
\paragraph{Mathematical
transformation}\label{mathematical-transformation}}

Because the detector cannot record new events while processing a
previous one, part of the acquisition time is effectively lost. This
effect is described by the \textbf{dead time}.

The dead time correction factor is computed from the ratio between real
acquisition time and live acquisition time:

\[
F_{\mathrm{dead}} = \frac{t_{\mathrm{real}}}{t_{\mathrm{live}}}
\]

where:

\begin{itemize}
\tightlist
\item
  \(t_{\mathrm{real}}\) is the total acquisition time,
\item
  \(t_{\mathrm{live}}\) is the effective counting time.
\end{itemize}

The corresponding dead-time percentage is:

\[
\mathrm{Dead\ Time}(\%) = \left(F_{\mathrm{dead}} - 1\right)\times 100
\]

This factor is written in the header of the exported \texttt{.txt} file
and is later reused during the ROI integration step to correct the
measured counts.

\hypertarget{error-sources}{%
\paragraph{Error sources}\label{error-sources}}

Several effects can affect the quality of the dead time estimate:

\begin{itemize}
\tightlist
\item
  Missing or malformed metadata in the original \texttt{.mpa} file.
\item
  Invalid or null live-time value, which prevents a meaningful
  correction factor from being computed.
\item
  Imperfect parsing if the raw acquisition file structure differs from
  the expected format.
\end{itemize}

If no valid value can be extracted, the correction cannot be reliably
applied and the associated spectrum must be interpreted with caution.

\hypertarget{impact-on-the-next-step}{%
\paragraph{Impact on the next step}\label{impact-on-the-next-step}}

The converted \texttt{.txt} files constitute the standard input format
for the rest of the application. They are used directly in the
calibration window and in the processing loop. {[}file:1{]}

The dead time factor stored in the file header is essential because it
is applied later to recover an estimate of the true number of detected
events inside the selected ROI. Any bias in this factor propagates to
the normalized counts and therefore to the excitation profile.

\begin{center}\rule{0.5\linewidth}{0.5pt}\end{center}

\hypertarget{energy-calibration}{%
\subsubsection{7.2 Energy Calibration}\label{energy-calibration}}

\hypertarget{input-data-1}{%
\paragraph{Input data}\label{input-data-1}}

The calibration step uses one converted spectrum \texttt{.txt} at a
time. Each file contains the ADC channel number and the corresponding
counts, together with the dead time factor written in the header during
the conversion stage.

The user identifies several known peaks in the spectrum and provides the
corresponding reference energies in keV. For each selected peak, the
software extracts a local fitting interval around the approximate peak
position.

\hypertarget{mathematical-transformation-1}{%
\paragraph{Mathematical
transformation}\label{mathematical-transformation-1}}

Each selected peak is fitted with a Gaussian-type model including a
linear background:

\[
f(x) = A \exp\left(-\frac{(x-x_0)^2}{2\sigma^2}\right) + b + cx
\]

where:

\begin{itemize}
\tightlist
\item
  \(A\) is the peak amplitude,
\item
  \(x_0\) is the centroid channel,
\item
  \(\sigma\) is the Gaussian width,
\item
  \(b\) and \(c\) describe the local linear background.
\end{itemize}

The fitted centroid \(x_0\) is taken as the best estimate of the channel
position of the calibration peak, and its uncertainty is extracted from
the covariance matrix of the fit.

Once several peaks have been identified, the software performs a
weighted linear regression between centroid channels and reference
energies according to:

\[
E = aC + b
\]

where:

\begin{itemize}
\tightlist
\item
  \(E\) is the energy in keV,
\item
  \(C\) is the channel number,
\item
  \(a\) is the calibration slope,
\item
  \(b\) is the calibration intercept.
\end{itemize}

The fit is weighted using the inverse squared uncertainty on the
centroid positions. The software also computes the uncertainties on
\(a\) and \(b\), as well as a quality indicator such as \(R^2\).

\hypertarget{error-sources-1}{%
\paragraph{Error sources}\label{error-sources-1}}

The calibration uncertainty is influenced by several effects:

\begin{itemize}
\tightlist
\item
  Inaccurate peak selection by the user.
\item
  Poor Gaussian convergence for weak, broad, or overlapping peaks.
\item
  Limited number of reference peaks.
\item
  Non-linearity of the detector response outside the fitted range.
\item
  Uncertainty on the reference energies themselves, if relevant in the
  experimental context.
\end{itemize}

If the selected peaks do not sufficiently constrain the calibration
line, the slope and intercept uncertainties increase and the energy
conversion becomes less reliable.

\hypertarget{impact-on-the-next-step-1}{%
\paragraph{Impact on the next step}\label{impact-on-the-next-step-1}}

The calibration coefficients \(a\) and \(b\) are written into the
configuration file and are later used to convert ROI limits expressed in
energy into channel limits used in the processing loop.

Any calibration error therefore propagates directly to the ROI
boundaries and affects the integrated counts computed in the next stage.

\begin{center}\rule{0.5\linewidth}{0.5pt}\end{center}

\hypertarget{roi-integration-and-normalization}{%
\subsubsection{7.3 ROI Integration and
Normalization}\label{roi-integration-and-normalization}}

\hypertarget{input-data-2}{%
\paragraph{Input data}\label{input-data-2}}

The ROI integration step uses:

\begin{itemize}
\tightlist
\item
  The converted spectral files generated during the conversion stage.
\item
  The calibration coefficients obtained during the calibration stage.
\item
  The ROI limits selected by the user in energy units.
\item
  The charge-related signal stored in the corresponding ADC channel used
  for normalization.
\item
  The experimental metadata stored in the configuration Excel file.
\end{itemize}

For each scenario, the program loops over a set of files and associates
them with the corresponding beam or terminal-voltage values extracted
from the measurement Excel files.

\hypertarget{mathematical-transformation-2}{%
\paragraph{Mathematical
transformation}\label{mathematical-transformation-2}}

The selected ROI is first converted from energy space to channel space
using the linear calibration:

\[
C = \frac{E-b}{a}
\]

where \(a\) and \(b\) are the slope and intercept determined during
calibration.

For each spectrum, the counts inside the ROI are then integrated. If the
ROI limits are not integer channel numbers, the boundary counts are
linearly interpolated so that fractional channel limits can be handled
consistently.

The total number of counts in the ROI is then corrected for dead time:

\[
N = N_{\mathrm{raw}} \times F_{\mathrm{dead}}
\]

The integrated charge is estimated from the dedicated charge channel and
scaled by a conversion factor used in the current implementation:

\[
Q = S_{\mathrm{charge}} \times 10^{-4}
\]

where \(S_{\mathrm{charge}}\) is the integrated value of the
charge-monitor channel.

The normalized count rate used to build the excitation profile is then:

\[
NC = \frac{N}{Q}
\]

This quantity is stored together with the corresponding beam energy for
each file in the loop.

\hypertarget{error-sources-2}{%
\paragraph{Error sources}\label{error-sources-2}}

Several factors contribute to the uncertainty on the normalized counts:

\begin{itemize}
\tightlist
\item
  Calibration uncertainty, which shifts the ROI boundaries.
\item
  Statistical uncertainty on the integrated counts, typically treated as
  Poissonian.
\item
  Uncertainty on the charge normalization.
\item
  Missing files or incomplete data in the acquisition sequence.
\item
  Interpolation errors at non-integer ROI boundaries.
\end{itemize}

These effects combine to determine the uncertainty associated with each
point of the excitation profile.

\hypertarget{impact-on-the-next-step-2}{%
\paragraph{Impact on the next step}\label{impact-on-the-next-step-2}}

The ROI integration and normalization step produces the excitation
profiles used in all subsequent analyses. Each output file contains
energy, normalized counts, and associated uncertainties.

These profiles are the direct input for optional peak removal and for
the final sigmoid fitting stage.

\begin{center}\rule{0.5\linewidth}{0.5pt}\end{center}

\hypertarget{peak-removal-and-data-cleaning}{%
\subsubsection{7.4 Peak Removal and Data
Cleaning}\label{peak-removal-and-data-cleaning}}

\hypertarget{input-data-3}{%
\paragraph{Input data}\label{input-data-3}}

The peak-removal step uses the excitation profiles generated by the loop
step, typically stored as Excel files with columns such as energy,
normalized counts, and uncertainties.

The user also defines a target energy, a search window around this
energy, and a half-width defining the interval to be removed when an
unwanted peak is detected.

\hypertarget{mathematical-transformation-3}{%
\paragraph{Mathematical
transformation}\label{mathematical-transformation-3}}

The program first restricts the profile to a local energy interval
around the target energy:

\[
[E_{\mathrm{center}} - \Delta E,\; E_{\mathrm{center}} + \Delta E]
\]

where \(E_{\mathrm{center}}\) is the selected resonance energy and
\(\Delta E\) is the search window.

Within this interval, a peak-detection algorithm is applied to the
normalized counts. If one or more local maxima are found, the detected
peak closest to the target energy is selected.

A removal interval is then defined around the detected peak energy
\(E_{\mathrm{peak}}\):

\[
[E_{\mathrm{peak}} - w,\; E_{\mathrm{peak}} + w]
\]

where \(w\) is the chosen half-width. All points located inside this
interval are removed from the profile, and the cleaned dataset is saved
as a new output file.

Diagnostic plots comparing raw and filtered profiles are also generated.

\hypertarget{error-sources-3}{%
\paragraph{Error sources}\label{error-sources-3}}

The peak-removal stage is sensitive to:

\begin{itemize}
\tightlist
\item
  False peak detection in noisy data.
\item
  Incorrect choice of the target energy or search window.
\item
  Excessive removal width, which may suppress useful physical
  information.
\item
  Insufficient removal width, which may leave part of the parasitic
  feature in the data.
\end{itemize}

Since this step modifies the dataset before fitting, it should be
applied with care and preferably validated visually using the generated
diagnostic plots.

\hypertarget{impact-on-the-next-step-3}{%
\paragraph{Impact on the next step}\label{impact-on-the-next-step-3}}

The cleaned profiles are stored separately and become the preferred
input for the sigmoid fit when a build-up peak would otherwise bias the
result.

If no peak is removed, the sigmoid fit can be applied directly to the
raw loop output instead.

\begin{center}\rule{0.5\linewidth}{0.5pt}\end{center}

\hypertarget{sigmoid-fitting-and-parameter-extraction}{%
\subsubsection{7.5 Sigmoid Fitting and Parameter
Extraction}\label{sigmoid-fitting-and-parameter-extraction}}

\hypertarget{input-data-4}{%
\paragraph{Input data}\label{input-data-4}}

The sigmoid fitting step takes as input either the raw excitation
profiles from the loop step or the cleaned profiles produced after peak
removal. These profiles contain energy, normalized counts, and
uncertainties.

Each profile is treated independently.

\hypertarget{mathematical-transformation-4}{%
\paragraph{Mathematical
transformation}\label{mathematical-transformation-4}}

The software fits each profile with a sigmoid model of the form:

\[
y(x) = \frac{L}{1+\exp\left[-k(x-x_0)\right]} + b
\]

where:

\begin{itemize}
\tightlist
\item
  \(L\) is the amplitude term,
\item
  \(x_0\) is the midpoint energy,
\item
  \(k\) is the steepness parameter,
\item
  \(b\) is the baseline offset.
\end{itemize}

From these parameters, the software derives a characteristic height
difference between the upper and lower plateaus:

\[
\Delta H = L - b
\]

The fitting routine also computes a goodness-of-fit indicator \(R^2\),
stores the fitted curve, and exports both numerical and graphical
outputs for each processed profile.

\hypertarget{error-sources-4}{%
\paragraph{Error sources}\label{error-sources-4}}

The fit quality depends on:

\begin{itemize}
\tightlist
\item
  The number and distribution of available points.
\item
  The presence of residual parasitic peaks or outliers.
\item
  The estimated uncertainties used as weights in the fit.
\item
  The suitability of the sigmoid model for the considered profile.
\end{itemize}

Profiles with too few valid points or strongly non-sigmoidal shapes may
lead to unstable or non-physical parameters.

\hypertarget{impact-on-the-next-step-4}{%
\paragraph{Impact on the next step}\label{impact-on-the-next-step-4}}

The parameters extracted from the sigmoid fit provide a compact
description of the excitation profile and can be used to compare
samples, group measurements, or perform higher-level statistical
analysis.

In the current workflow, the plateau difference and related fit
parameters are among the main quantities retained for interpretation.

\begin{center}\rule{0.5\linewidth}{0.5pt}\end{center}

\hypertarget{uncertainty-estimation-at-each-step}{%
\subsubsection{7.6 Uncertainty Estimation at Each
Step}\label{uncertainty-estimation-at-each-step}}

\hypertarget{input-data-5}{%
\paragraph{Input data}\label{input-data-5}}

Uncertainty propagation combines information from several stages of the
workflow:

\begin{itemize}
\tightlist
\item
  The uncertainty on peak centroid positions from Gaussian fitting.
\item
  The uncertainty on calibration coefficients.
\item
  The statistical uncertainty on integrated counts.
\item
  The uncertainty associated with charge normalization.
\item
  The uncertainties stored in the processed excitation profiles.
\end{itemize}

\hypertarget{mathematical-transformation-5}{%
\paragraph{Mathematical
transformation}\label{mathematical-transformation-5}}

At the calibration stage, the uncertainty on each centroid is extracted
from the covariance matrix of the Gaussian fit and used as the weight in
the linear energy calibration.

During ROI integration, the code combines a counting-statistics term and
a calibration-related term for the corrected counts \(N\). In the
current implementation, the statistical term is approximated by:

\[
\sigma_{\mathrm{stat},N} = \sqrt{N}
\]

and the calibration-related contribution is written as:

\[
\sigma_{\mathrm{cal},N} = \frac{\varepsilon_{\mathrm{cal}}}{100}\,N
\]

where \(\varepsilon_{\mathrm{cal}}\) is the relative calibration error
in percent.

The total uncertainty on \(N\) is then combined quadratically:

\[
\sigma_N = \sqrt{\sigma_{\mathrm{stat},N}^2 + \sigma_{\mathrm{cal},N}^2}
\]

An uncertainty is also associated with the normalized charge term, and
the uncertainty on the normalized count \(NC=N/Q\) is propagated using
the usual relative-error formula:

\[
\sigma_{NC} = NC \sqrt{\left(\frac{\sigma_N}{N}\right)^2 + \left(\frac{\sigma_Q}{Q}\right)^2}
\]

These values are stored in the output profiles and may be used later as
weights in the sigmoid fit.

\hypertarget{error-sources-5}{%
\paragraph{Error sources}\label{error-sources-5}}

The uncertainty model currently implemented includes the main
experimental contributions, but it remains an approximation. Its
accuracy depends on:

\begin{itemize}
\tightlist
\item
  The validity of the Poisson approximation for the corrected counts.
\item
  The quality of the calibration error estimate.
\item
  The reliability of the charge normalization factor.
\item
  The assumption that different sources of uncertainty are independent.
\end{itemize}

Additional systematic effects may still need to be incorporated
depending on the final scientific use of the results.

\hypertarget{impact-on-the-next-step-5}{%
\paragraph{Impact on the next step}\label{impact-on-the-next-step-5}}

Uncertainty estimates are essential for judging the quality and
comparability of the processed profiles. They also affect the weighting
of the sigmoid fit and therefore the stability of the extracted
parameters.

A realistic uncertainty budget is therefore required before any robust
statistical comparison between samples can be performed. \#\#\# 7.7
Statistical Comparison Using ANOVA

\hypertarget{code-architecture-optional}{%
\subsection{8. Code Architecture
(optional)}\label{code-architecture-optional}}

\hypertarget{limitations}{%
\subsection{9. Limitations}\label{limitations}}

\hypertarget{references}{%
\subsection{10. References}\label{references}}
